% !TeX root = main.tex
\lecture{20}{Thu 05 Mar 2026 09:00}{Methods of Integration}

\begin{tcolorbox}
    The derivative of an elementary function is always an elementary function. The integral of an elementary function is not always elementary however.

    As a result, differentiation is always purely algorithmic churn, integration is often tricky.
\end{tcolorbox}

Methods of integration are often based on methods of differnetiation, for example:
\begin{itemize}
    \item Integration by inspection: just spotting the answer.
    \item Integration by parts: Inverse of the product rule of differentiation.
    \item Integration by substitution: Inverse of the chain rule of differentiation.
\end{itemize}

\section{Integration By Inspection}
This is the most difficult to codify and can be easy for some people and very easy for others.

This might involve:
\begin{itemize}
    \item Trigonometric Identities, for example replacing an expression whose integral is unknown with an expression whose integral is, using a trig sub.
\end{itemize}

